\section*{Chapitre \ref{ch:rc:commodity-and-energy-derivatives} --- Dérivés sur matières premières et sur énergie}

\begin{solution}{exo:rc:commodity-and-energy-derivatives:1}
27.2~\%, contre 68.7~\% à la livraison. Le poids du facteur de court terme, $e^{-\kappa\tau}$, a
décru jusqu'à $e^{-1.5} = 0.22$~: une livraison à un an évolue surtout avec le niveau de long terme,
car les chocs sur l'offre et la demande d'aujourd'hui devraient s'estomper avant cette livraison.
\end{solution}

\begin{solution}{exo:rc:commodity-and-energy-derivatives:2}
Tout programme achète et vend au même prix à terme et supporte les frais d'injection et
de soutirage, si bien que le meilleur plan déterministe consiste à ne rien faire. L'installation
reste néanmoins un portefeuille d'\omterm{def:rc:commodity-and-energy-derivatives:spread}{options sur spread}~: à mesure que les prix bougent, les écarts entre mois
s'ouvrent et son propriétaire peut les exploiter~; sa valeur est entièrement extrinsèque.
\end{solution}

\begin{solution}{exo:rc:commodity-and-energy-derivatives:3}
Sa valeur intrinsèque, actualisée d'un jour~: $3.50-2.74-0.50 = 0.26$ dollars par MMBtu
(0.2600 après actualisation d'un jour à 4~\%).
\end{solution}

\begin{solution}{exo:rc:commodity-and-energy-derivatives:4}
La volatilité du spread vaut $\sqrt{\sigma_1^2+\sigma_2^2-2\rho\sigma_1\sigma_2}$ (à strike nul)~: une corrélation plus forte
signifie un spread moins volatil et une option moins chère. Les mois de livraison proches sont fortement
corrélés (0.970 pour juillet et janvier à l'échéance de juillet dans le chapitre), les mois éloignés moins.
\end{solution}

\begin{solution}{exo:rc:commodity-and-energy-derivatives:5}
Il part du plan intrinsèque couvert par des contrats à terme, et ne modifie le plan que lorsque le
nouveau plan, évalué sur la nouvelle courbe, vaut davantage que l'ancien une fois la couverture
renouvelée~: chaque modification ajoute une valeur positive ou nulle, figée par les contrats à terme.
\end{solution}

\begin{solution}{exo:rc:commodity-and-energy-derivatives:6}
Le strip donne le droit d'exercer chaque call mensuel~; le swing plafonne le total, si bien que son
détenteur doit choisir les mois à utiliser et renonce ainsi à certains exercices que le strip prendrait.
Le strip est une borne supérieure (6.11 contre 4.37 dans le chapitre).
\end{solution}

\begin{solution}{exo:rc:commodity-and-energy-derivatives:7}
À strike nul, la formule de Kirk se confond avec celle de Margrabe, exacte pour des prix à terme lognormaux~: la différence
avec la méthode de Monte-Carlo est de 0.007 cent, dans le bruit. Pour un strike de 0.50, il vaut $-0.375$ cent~: la somme
$F_2+K$ n'est pas lognormale, et l'approximation de sa volatilité par $\sigma_2F_2/(F_2+K)$ est
peu fiable au voisinage de la monnaie, là où l'option est la plus sensible à la forme de la distribution.
\end{solution}

\begin{solution}{exo:rc:commodity-and-energy-derivatives:8}
Un book couvert à l'intrinsèque ne l'est que contre les déplacements parallèles de la courbe à terme, mais il
détient des options~: en février 2021, le droit de soutirer rapidement au comptant valait des millions de
dollars pour une installation rapide, et un book qui ne l'a pas exercé y a renoncé. À l'inverse, un
book vendeur de flexibilité (contrats d'approvisionnement avec swing) a subi une lourde perte. La couverture
intrinsèque ne mesure ni l'un ni l'autre.
\end{solution}

\begin{solution}{pb:rc:commodity-and-energy-derivatives:1}
\textbf{1.} USD~634\,447, 699\,550 et 698\,182.
\textbf{2.} Environ USD~64\,000, soit 10~\% de l'intrinsèque~: le droit de modifier le programme à mesure que
les écarts entre mois évoluent.
\textbf{3.} La volatilité de court terme $\sigma_s$, avec le retour à la moyenne $\kappa$~: ils fixent dans quelle mesure
les mois proches évoluent l'un contre l'autre.
\textbf{4.} L'installation n'agit qu'une fois par mois sur des prix mensuels~; une caverne de sel agit chaque jour,
sur des prix quotidiens dont les pics se compensent dans un prix mensuel.
\textbf{5.} Acheter les contrats à terme d'été et vendre ceux d'hiver du plan intrinsèque, aux
volumes du plan.
\textbf{6.} 2.88 dollars le 1 février~; 23.86 le 17 février.
\textbf{7.} 2.62 dollars.
\textbf{8.} 500\,000 MMBtu pendant cinq jours (11, 12, 16, 17 et 18 février), soit un gain de
USD~4\,306\,870.
\textbf{9.} 90\,000 MMBtu pendant neuf jours, soit un gain de USD~479\,057.
\textbf{10.} La caverne épuise son gaz au bout de cinq jours et ne vend que les meilleurs~; le
réservoir vend un peu chaque jour, y compris les jours où le prix va de 3.35 à 3.76 dollars.
\textbf{11.} $\ln(23.86/3.76) = 1.85$ contre $0.687\sqrt{4/252} = 0.087$~: 21 écarts-types,
impossible dans le modèle.
\textbf{12.} Des sauts (un spot en diffusion à sauts ou à changement de régime), ou une volatilité stochastique
qui augmente avec la rareté.
\textbf{13.} Le pic a duré une semaine et a laissé la courbe à terme de la saison suivante presque
inchangée, si bien que le spread saisonnier (la valeur intrinsèque) n'a pas bougé~; les options ont gagné
parce qu'elles ont été exercées pendant le pic.
\textbf{14.} Les contrats swing et autres droits de prise de gaz à prix fixe, la capacité ferme de transport
hors de la région gelée, et des options d'achat sur gaz au jour le jour.
\textbf{15.} Les vendeurs de ces droits~: les fournisseurs soumis à des obligations swing, qui devaient acheter du gaz
au comptant pour livrer au prix du contrat, et les producteurs d'électricité à découvert de combustible pour leurs ventes d'électricité.
\textbf{16.} La valeur se concentre dans des semaines rares et extrêmes~; l'année moyenne la sous-estime, et un
seul événement peut couvrir des années de frais.
\textbf{17.} La position intrinsèque et ses couvertures, la \omterm{def:rc:commodity-and-energy-derivatives:extrinsic}{valeur extrinsèque} et sa sensibilité à
la volatilité, la valeur dans les semaines de stress historiques, et les contraintes physiques (cadences,
ratchets, contrats susceptibles d'imposer des livraisons).
\textbf{18.} Ajouter des pics (sauts ou régimes) au facteur de court terme, et décider au jour le jour.
\textbf{19.} Cas nommé~: \emph{la caverne de sel en février 2021}~: l'installation rapide gagne
USD~4\,306\,870 sur 500\,000 MMBtu, la lente USD~479\,057 sur 90\,000.
\textbf{20.} Les options permettant de déplacer du gaz dans le temps à court préavis, qui rapportent le plus
précisément dans les semaines qu'aucun modèle de la moyenne ne prédit.
\end{solution}

\begin{solution}{iq:rc:commodity-and-energy-derivatives:1}
Les contrats à courte échéance sont plus volatils que ceux à longue échéance. Dans un modèle à deux facteurs, un
facteur de court terme à retour à la moyenne agit sur les contrats proches avec le poids $e^{-\kappa\tau}$, qui s'estompe
pour les livraisons lointaines, tandis qu'un facteur de long terme les fait tous bouger.
\iqlookfor{l'observation et le mécanisme.}
\end{solution}

\begin{solution}{iq:rc:commodity-and-energy-derivatives:2}
Strike nul~: Margrabe, exact pour des prix lognormaux. Strike non nul~: \omterm{def:rc:commodity-and-energy-derivatives:kirk}{approximation de Kirk},
le raffinement de Bjerksund--Stensland, l'intégration unidimensionnelle conditionnée à un prix,
ou la méthode de Monte-Carlo. L'\omterm{def:rc:commodity-and-energy-derivatives:kirk}{approximation de Kirk} est précise pour de petits strikes par rapport au second prix, et moins
au voisinage de la monnaie, ou pour de grands strikes et des volatilités très différentes.
\iqlookfor{Margrabe, Kirk et quand intégrer.}
\end{solution}

\begin{solution}{iq:rc:commodity-and-energy-derivatives:3}
Intrinsèque~: optimiser le programme sur les contrats à terme du jour et le couvrir. \omterm{def:rc:commodity-and-energy-derivatives:extrinsic}{Intrinsèque glissant} (rolling intrinsic)~:
réoptimiser et réajuster la couverture à mesure que la courbe évolue, borne inférieure de l'optionnalité totale. Complète~: une
programmation dynamique stochastique sur le spot (moindres carrés de Monte-Carlo ou arbre), évaluant les
décisions quotidiennes sous contraintes.
\iqlookfor{les trois couches et leur ordre.}
\end{solution}

\begin{solution}{iq:rc:commodity-and-energy-derivatives:4}
Simuler le spot~; induction rétrograde sur une grille d'inventaire~; par période et à l'inventaire final,
régresser la valeur future réalisée sur des fonctions de base de l'état~; agir sur les estimations, conserver
les valeurs réalisées. Biais~: anticipation en échantillon (élevé), politiques de régression sous-optimales (faible)~;
utiliser des trajectoires indépendantes pour la passe d'évaluation et la comparer à l'\omterm{def:rc:commodity-and-energy-derivatives:extrinsic}{intrinsèque glissant}.
\iqlookfor{l'algorithme et les deux biais.}
\end{solution}

\begin{solution}{iq:rc:commodity-and-energy-derivatives:5}
Couvrir l'intrinsèque par des contrats à terme et réajuster la couverture à mesure que le plan change~; mesurer
la \omterm{def:rc:commodity-and-energy-derivatives:extrinsic}{valeur extrinsèque} et sa sensibilité à la volatilité~; tester sous stress avec des pics historiques et des semaines froides~; suivre
les contraintes physiques et l'inventaire~; limiter la flexibilité vendue à découvert dans les contrats swing~;
conserver de la liquidité pour les appels de marge sur les couvertures.
\iqlookfor{couverture, stress et contraintes physiques.}
\end{solution}

\begin{solution}{iq:rc:commodity-and-energy-derivatives:6}
Des facteurs gaussiens donnent des prix lognormaux à queues fines~; les pics viennent d'une offre et d'une demande
inélastiques et durent quelques jours. Ajouter des sauts à retour rapide à la moyenne, un changement de régime, ou un
modèle fondamental d'offre et de demande avec un empilement de coûts très pentu.
\iqlookfor{queues fines et remède à retour rapide.}
\end{solution}
